% Isolated post-CP28 candidate. No exclusion of the larger wild images.
\subsection{The minimal third-level wild image at the first wild index}

Throughout, \(p\ge5\), \(q=p^2\), \(N=q-1\), and \(K\) is a
complete discretely valued field of mixed characteristic \((0,p)\),
with algebraically closed residue field, \(\mu_p\subset K\), and
absolute index \(e=p(p-1)\). All torsion modules below are actual
generic torsion of height-two groups with potentially supersingular
good reduction. Their finite special characters use the covariant
convention fixed in B53.

\begin{lemma}[Third-level valuations and the two wild layers]
\label{lem:B-RD79-layers}
\status{Candidate; complete proof supplied, independent audits pending.}
Suppose the residual inertia module of \(M\) is irreducible over
\(\mathbf F_p\). Let \(L/K\) be the good tame extension of degree
\(N\), and put \(F=L(M[p^3])\), \(W=\operatorname{Gal}(F/L)\).
Then \(M[p]\) is \(L\)-rational, and every point of exact order
\(p^j\), \(1\le j\le3\), has coordinate valuation
\[
             v_L(T(P))=\frac{e}{q^{j-1}}.
                                                     \tag{B-RD79.1}
\]
Consequently \(W\) has order \(p^r\), where \(3\le r\le8\).
If \(r=3\), its reduction image and terminal kernel are
\[
 \begin{split}
 W_1&=\{A\bmod p:1+pA+p^2B\in W\}\subset M_2(\mathbf F_p),\\
 W_2&=\{B\bmod p:1+p^2B\in W\}\subset M_2(\mathbf F_p),\\
 \dim_{\mathbf F_p}W_1&=1,\qquad
 \dim_{\mathbf F_p}W_2=2,\qquad W_1\subset W_2.
 \end{split}                                         \tag{B-RD79.2}
\]
The inclusion in the last line is in the same actual matrix
coordinates, induced by taking \(p\)th powers.
\end{lemma}
\begin{proof}
A proper nonzero valuation-filtration step in \(M[p]\) would be
inertia stable, contradicting irreducibility. Thus all its nonzero
points have one valuation. The product of the \(q-1\) nonzero
parameters has valuation \(v_L(p)=eN\), so this valuation is \(e\).
The normal wild subgroup fixes a nonzero vector in characteristic
\(p\); its fixed space is inertia stable. Irreducibility therefore
makes wild inertia trivial. The remaining cyclic irreducible
image embeds in \(\mathbf F_q^\times\), and is killed by the tame
extension of degree \(N\). The good-field assertion itself is
Lemma~\ref{lem:B-RD49-inertia}, which permits arbitrary \(e\).

The first-level Newton polygon is the segment from \((1,eN)\)
to \((q,0)\). Its intermediate degree-\(i\) coefficients have
valuation at least \((q-i)e\). For a target of valuation
\(0<Z\le e\), these coefficients lie strictly above the segment
from \((0,Z)\) to \((q,0)\): for \(0<i<q\),
\((q-i)e>(1-i/q)Z\). Hence all its \(q\) inverse points have
valuation \(Z/q\). Applying this twice proves (B-RD79.1).
This is the same exact inverse-polygon argument as
Lemma~\ref{lem:B-RD55-valuations}; no approximation by an abstract
formal module is made.

Since the residual module is rational over \(L\), the group \(W\)
lies in the kernel of
\(\GL_2(\mathbf Z/p^3)\to\GL_2(\mathbf F_p)\), of order \(p^8\).
All residue extensions are trivial. Every primitive third-level
point has valuation
\(e/q^2=(p-1)/p^3\), whose reduced denominator is \(p^3\).
Its point field therefore has degree at least \(p^3\) over \(L\).
This proves \(3\le r\le8\).

Both spaces in (B-RD79.2) are additive, and
\(r=\dim W_1+\dim W_2\). For \(p\ge3\), the exact matrix identity
\[
       (1+pA+p^2B)^p=1+p^2A\pmod{p^3}
                                                     \tag{B-RD79.3}
\]
shows \(W_1\subset W_2\). If \(W_1=0\), each primitive point
has at most \(p^2\) conjugates: a terminal-layer element moves
it by an element of its first-level module, of cardinality \(p^2\).
This contradicts its valuation denominator. Thus, when \(r=3\),
the nonnegative dimensions with sum three, with the first positive
and no greater than the second, are exactly \((1,2)\).
\end{proof}

\begin{proposition}[The minimal third-level image cannot conceal the character]
\label{prop:B-RD79-minimal}
\status{Candidate; complete proof supplied, independent audits pending.}
Let \(M_1[p^3]\simeq M_2[p^3]\) be an actual isomorphism with
common irreducible residual inertia module. If its wild image
has order \(p^3\), the two finite special characters agree.
There is no CM assumption.
\end{proposition}
\begin{proof}
Use the same actual good field \(L/K\) of degree \(N\) and
common full division field
\(F=L(M_1[p^3])=L(M_2[p^3])\). Put \(R_i=M_i[p^3](F)\) and
\(V_i=R_i/pR_i\). The preceding lemma and \([F:L]=p^3\) give
\[
 \begin{gathered}
 A=v_F(T_i(P))=p-1\quad(P\text{ primitive of order }p^3),\\
 v_F(T_i(Q))=qA\quad(Q\text{ of exact order }p^2),\qquad
 v_F(T_i(S))=q^2A\quad(S\ne0\text{ of order }p),\\
                           v_F(p)=Nq^2A.
 \end{gathered}                                        \tag{B-RD79.4}
\]
Choose a common uniformizer \(\varpi\) of \(F\) and special
Honda coordinates with \([p](T)\equiv T^q\) on the special
fibre. Leading coefficients give injective linear maps
\[
       \ell_i:V_i\longrightarrow k,\qquad
       \ell_i(P)=\overline{T_i(P)/\varpi^A}.
                                                     \tag{B-RD79.5}
\]
Indeed two primitive points have the same leading coefficient
exactly when their difference lies in \(pR_i\): the latter
points have valuation at least \(qA>A\), whereas every point
outside \(pR_i\) has valuation \(A\). Integral formal addition
is ordinary addition in the leading term. The inverse-polygon
calculation shows that the degree-\(q\) term is uniquely lowest
at every primitive third- or second-level point. Consequently
multiplication by \(p^2\) sends leading coordinates to
\(\ell_i(P)^{q^2}\).

The first-level leading space is \(d\mathbf F_q\) for a nonzero
\(d\). To see this with the present absolute index, let
\(\zeta\) be the chosen primitive tame \(N\)th root. If the
special tangent exponent is \(c_i\), the first-level leading
action is multiplication by \(\zeta^{c_i+e}\), using
\(v_L(T(S))=e\). It is not in \(\mathbf F_p\) by residual
irreducibility. A two-dimensional \(\mathbf F_p\)-space stable
under this scalar is a one-dimensional \(\mathbf F_q\)-space.
Taking the unique \(q^2\)th roots in the algebraically closed
residue field shows that the image of (B-RD79.5) is likewise
\(d_i\mathbf F_q\), for some \(d_i\ne0\). Normalize it by
\(\ell_i(v)=d_i z\); its first-level normalization is then
\(d_i^{q^2}z\), since \(z^{q^2}=z\) on \(\mathbf F_q\).

Write \(U=\ker(W\to W_1)\), so that \(|U|=q\). For a
nonidentity \(g\in U\), its action on \(R_i\) is
\(1+p^2 B_i(g)\). Let \(b_g\) be its actual lower break and
\[
 d_g=\overline{(g\varpi-\varpi)/\varpi^{b_g+1}}\ne0.
\]
Choose residue representatives in a coefficient ring inside
\(\mathcal O_K\), as in the complete fixed-coefficient proof
of Lemma~\ref{lem:B-RD60-break} and
\cite[Theorem~10.160.8 and proof]{Stacks032A}.
Every \(g\) fixes these representatives. Since \(p\nmid A\),
the same uniformizer expansion proves for every \(x\) of
valuation \(A\)
\[
 v_F(gx-x)=A+b_g,\qquad
 \overline{(gx-x)/\varpi^{A+b_g}}
       =\overline A\,d_g\,\overline{x/\varpi^A}.
                                                     \tag{B-RD79.6}
\]
Explicitly the first binomial change of its degree-\(A\) term
is nonzero; every higher degree changes in valuation at least
its degree plus \(b_g\), strictly later. This accounts for
all unit terms, rather than replacing a point by a monomial.

Thus \(g\) moves every primitive point. Its displacement is a
nonzero first-level point, so (B-RD79.4) and (B-RD79.6) force
\[
          b_g=(q^2-1)A\quad(g\in U\setminus\{1\}).
                                                     \tag{B-RD79.7}
\]
Formal and ordinary coordinate subtraction have the same
leading coefficient. Applying the \(p^2\)-series leading term
therefore gives
\[
   \ell_i(B_i(g)v)^{q^2}
               =\overline A\,d_g\,\ell_i(v).
                                                     \tag{B-RD79.8}
\]
In the normalized field coordinates this says that \(B_i(g)\)
is multiplication by
\[
        \lambda_i(g)=\kappa_i d_g\in\mathbf F_q,
             \qquad \kappa_i=\overline A d_i^{1-q^2}.
                                                     \tag{B-RD79.9}
\]
The representation of \(U\) by these matrices is injective,
so \(\lambda_i(U)=\mathbf F_q\), with zero assigned to the
identity. In particular this is the full Cartan plane, obtained
from the actual last-layer action, not assumed from its dimension.

Suppose the requested finite characters differed. The isomorphism
on the first-level leading spaces must be of the form
\(z\mapsto a z^p\), \(a\in\mathbf F_q^\times\).
Here is the exact reason, independent of the index restriction
in the earlier B53 statement. The common irreducible tame
representations have eigenvalue orbits
\(\{\zeta^{c_i+e},\zeta^{p(c_i+e)}\}\).
An intertwiner is either field-linear, identifying \(c_1=c_2\)
modulo \(N\), or Frobenius-semilinear. The field-linear case
would give equal special characters, contrary to assumption.
The isomorphism on \(V_i\) is the same one under multiplication
by \(p^2\), since that multiplication raises normalized leading
coordinates to \(z^{q^2}=z\).
It follows that
\[
            \lambda_2(g)=\lambda_1(g)^p
                        \quad\text{for every }g\in U.
                                                     \tag{B-RD79.10}
\]
But (B-RD79.9) simultaneously makes \(\lambda_2(g)\) a fixed
scalar multiple of \(\lambda_1(g)\). Since the latter run through
all of \(\mathbf F_q\), this would make its nontrivial Frobenius
a scalar map. Evaluation at one and at an element outside
\(\mathbf F_p\) is a contradiction. Hence the special characters
agree.
\end{proof}

\begin{corollary}[A stricter remaining irreducible-residual gate]
\label{cor:B-RD79}
\status{Candidate; complete proof supplied, independent audits pending.}
At the first wild absolute index \(e=p(p-1)\), any actual
\(p^3\)-congruent pair with irreducible residual inertia and
different finite special characters has wild image of order
\(p^r\) with \(4\le r\le8\).
In particular, at \(p=5,e=20\), its common actual \(125\)-torsion
wild image has order at least \(625\).
\end{corollary}
\begin{proof}
The first lemma gives \(3\le r\le8\), and the proposition
excludes \(r=3\). The conclusion concerns the actual image,
not a conjectural formal-group model with that image size.
\end{proof}

\begin{remark}[Scope of the additional image restriction]
\label{rem:B-RD79-scope}
\status{Candidate scope; larger images and reducible residual cases remain open.}
This argument imposes no potential-CM hypothesis. It uses the
terminal kernel of an actual third-level representation and
the actual common division field; it does not identify that
kernel with the entire wild image. Neither any of the orders
\(p^4,\ldots,p^8\), nor the reducible-residual comparisons, is
excluded. In particular the unrestricted first-wild least
index remains in \(\{2,3\}\). The A/C finite-type relation alone
is not treated as an actual quadratic self-twist.


\smallskip
\noindent\emph{Exact supplement and provenance.}
The retained standalone program
\path{code/reduction_descent/verify_first_wild_irreducible_image.py}
checks 216 valuation-denominator cases, 600 layer-dimension
cases, 24 all-binomial-coefficient divisibility rows, 960
higher uniformizer terms, 75,984 \(q^2\)-power field identities,
75,984 distinct scalar candidates for Frobenius, and 88,335
literal matrix \(p\)th-power identities modulo \(p^3\).
The author run passed in \(0.9368512359942542\) seconds.
It does not construct a larger image or a formal-group pair.
No executed computation failed. The full fixed-coefficient
uniformizer and leading-coordinate arguments of B55, and the
actual inverse-polygon argument, were reread rather than
importing a bound from an abstract group order. No new
external-source theorem is asserted; the already freshly read
Stacks coefficient-ring proof supplies the fixed representatives.
\end{remark}
